
\svnid{$Id: $}
\svnid{$URL: $}

\section{FOR REVIEW:1D universal weir for steady flow conditions}
\newrefsegment

%---------------------------------------------------------------------------------
\section*{Quality Assurance}
\begin{tabular}{@{}p{27mm}@{}p{0mm}p{\textwidth-52mm-48pt}p{15mm}p{10mm}}
\textbf{Date} & \textbf{:} & October 2019   \\ 
\textbf{Author} & \textbf{:} & Eskedar Gebremedhin  & \textbf{Initials:} & E.T. \\
\textbf{Review} & \textbf{:} &   & \textbf{Initials:} & \ldots 
\end{tabular}

%---------------------------------------------------------------------------------
\section*{Version information}
{{
\begin{tabular}{@{}p{27mm}@{}p{0mm}p{\textwidth-27mm-24pt}}
\textbf{Executable}   & \textbf{:} & Deltares, D-Flow FM Version 1.2.67.64909M, October 2019, 10:43:09  \\
\textbf{Location input}     & \textbf{:} & \url{https://repos.deltares.nl/repos/DSCTestbench/trunk/cases/e02_dflowfm/f106_hydraulic-structures/c06_Universal_Weir} \\
\textbf{Revision input} & \textbf{:} & - \\
\textbf{Location output}     & \textbf{:} & \url{https://repos.deltares.nl/repos/DSCTestbench/trunk/references/win64/e02_dflowfm/f106_hydraulic-structures/c06_Universal_Weir} \\
\textbf{Revision output} & \textbf{:} & -
\end{tabular}
}}

%---------------------------------------------------------------------------------
\section*{Purpose}
The purpose of this validation test case is to verify that the structure equations of 1D universal weir for steady flow conditions are solved correctly in \dflowfmfull and that the results produced are similar to the analytical solution and those of \sobekthree model.

%---------------------------------------------------------------------------------
\section*{Linked claims}
\begin{enumerate}
\item The structure equations of a 1D universal weir are solved correctly in \dflowfmfull.
\item The modelled discharges through the universal weir are similar to the analytical solution and \sobekthree model result.
\end{enumerate}

%---------------------------------------------------------------------------------
\section*{Approach}
To verify the claims, the modelled discharges of two straight channels with opposite flow direction for steady-state are compared to the analytical solution calculated using the universal weir equations (see \dflowfm technical user manual for the equations). 
Furthermore, the modelled discharges in \dflowfm are compared with \sobekthree model output.

%---------------------------------------------------------------------------------
\section*{Conclusion} 
The \dflowfm results match the analytical solution. Additionally, the \dflowfm result is very similar to the \sobekthree model result.
It is therefore concluded that the 1D universal weir equations are solved correctly in \dflowfm.

%---------------------------------------------------------------------------------
\section*{Model description}
The test comprises of two separate sub-models, T1 and T2 (see \autoref{fig:e02_f106_c06_Test_model_lay-out}). 
T1 flows from West to East and T2 flows from East to West.
The sub-models are 110 meters in length and have an equidistant grid spacing of 10m. 
The boundary conditions, cross-sectional profiles and hydraulic roughness of the sub-models are given on \autoref{tab:e02_f106_c06_Boundary_conditions} and \autoref{tab:e02_f106_c06_Cross_Sections}.
Universal weir structures are placed at the centre of the sub-models at chainage=55m.
\autoref{tab:e02_f106_c06_Structure} and \autoref{tab:e02_f106_c06_weir_shape} lists the weir parameters and shape used in this test case.

\begin{figure}[H]
    \centering
    \includegraphics*[width=0.75\textwidth]{figures/e02_f106_c06_layout.JPG}
    \caption{Lay-out of the test sub-models}
	\label{fig:e02_f106_c06_Test_model_lay-out}
\end{figure}

\begin{longtable}{|p{18mm}|p{23mm}|p{35.0mm}|p{65mm}|}
\caption{Boundary conditions of the test sub-models} 
\label{tab:e02_f106_c06_Boundary_conditions} \\ 
\hline
\STRUT
      {Sub-model} 
	 &  {Chainage \unitbrackets{m}}  
	 &  {ID}  
	 &  {Boundary condition}  \\ 
[1ex]\hline \endhead
\hline
\endfoot
%
\STRUT										   
{T1} & {0.00}    & {T1\_B1\_WE\_x0m}     &  {Time Varying Water Levels on \autoref{fig:e02_f106_c06_Upstream_Water_Level_Boundary}}   \\
                                               & {110.00}        & {T1\_B1\_WE\_x110m}         & {Time Varying Water Levels on \autoref{fig:e02_f106_c06_Downstream_Water_Level_Boundary}}                     \\
\hline
\STRUT											   
{T2} & {0.00}     & {T2\_B1\_EW\_x0m}  & {Time Varying Water Levels on \autoref{fig:e02_f106_c06_Upstream_Water_Level_Boundary}}   \\
                                               & {110.00}        & {T2\_B1\_EW\_x110m}      & {Time Varying Water Levels on \autoref{fig:e02_f106_c06_Downstream_Water_Level_Boundary}}                     \\ [1ex]
                                                                                                                              
\end{longtable}	 


\begin{figure}[H]
    \centering
    \includegraphics*[width=0.95\textwidth]{figures/e02_f106_c06_Upstream_Boundary.JPG}
    \caption{Upstream Water Level Boundary}
		\label{fig:e02_f106_c06_Upstream_Water_Level_Boundary}
\end{figure}

\begin{figure}[H]
    \centering
    \includegraphics*[width=0.95\textwidth]{figures/e02_f106_c06_Downstream_Boundary.JPG}
    \caption{Downstream Water Level Boundary}
		\label{fig:e02_f106_c06_Downstream_Water_Level_Boundary}
\end{figure}

\begin{longtable}{|p{10mm}|p{15mm}|p{28mm}|p{10mm}|p{10mm}|p{18mm}|p{25mm}|}
\caption{ZW Cross Section Definition of the test sub-models} 
\label{tab:e02_f106_c06_Cross_Sections} \\ 
\hline
\STRUT
      {Sub-model} 
	 &  {Chainage \unitbrackets{m}}  
	 &  {Cross-Section ID}  
	 &  {Width \unitbrackets{m}}  
	 &  {Height \unitbrackets{m}} 
	 &  {Bed Level \unitbrackets{mAD}} 
	 &  {Chezy Roughness \unitbrackets{m\textsuperscript{0.5}/s}}\\ 
[1ex]\hline \endhead
\hline
\endfoot
%
\STRUT										   
{T1} & {0.00}    & {T1\_B1\_cross\_x0}     &  {50.00} &  {50.00}  &  {0.00} &  {50.00}     \\
\STRUT											   
{T2} & {0.00}    & {T2\_B1\_cross\_x0}     &  {50.00} &  {50.00}  &  {0.00} &  {50.00}     \\ [1ex]                                                                                                                       
\end{longtable}	 


\begin{longtable}{|p{58.0mm}|p{29.0mm}|p{29.0mm}|}
	\caption [Structure characteristics] {Characteristics of structure} 
	\label{tab:e02_f106_c06_Structure} \\ 
	\hline
	\STRUT
	    {Parameter} 
	  & {T1}
	  & {T2}  \\
	\hline \endhead
	\hline
	%
	\STRUT
	{Structure ID}                                       & {T1\_B1\_weir\_x55m} 	& {T2\_B1\_weir\_x55m}  \\	              
	\STRUT
	{Chainage in x-direction \unitbrackets{$m$}}         & {55.0} 				    & {55.0}				\\
	\STRUT
	{Crest level \unitbrackets{$m AD$}}                  & {2.00}                   & {2.00}                \\
	\STRUT
	{Crest width \unitbrackets{$m$}}  			         & {10.00}                  & {10.00}               \\
	\STRUT
	{allowedFlowDir}                                     & {both}                   & {both}               \\
	\STRUT
	{corrCoeff}                                          & {1}                    & {1}                \\
	[1ex] \hline
\end{longtable}


\begin{longtable}{|p{29.0mm}|p{29.0mm}|}
	\caption [Asymmetric shape of the universal weirs] {Asymmetric shape of the universal weirs} 
	\label{tab:e02_f106_c06_weir_shape} \\ 
	\hline
	\STRUT
	    {Y}
	  & {Z}  \\
	\hline \endhead
	\hline
	%
	\STRUT
	{0.00} 	& {2.50}  \\	              
	\STRUT
	{2.49} 	& {2.50}  \\
	\STRUT				
	{2.50}  & {2.00}  \\
	\STRUT				
	{7.50}  & {2.00}  \\
	\STRUT				
	{7.51}  & {2.90}  \\
	\STRUT				
	{10.00} & {2.90}  \\
	[1ex] \hline
\end{longtable}

%---------------------------------------------------------------------------------
\section*{Results}
The modelled and calculated discharges through the universal weir for sub-model T1 and T2 are similar. The results are presented on \autoref{tab:e02_f106_c06_universal_weir_T1_T2} and \autoref{fig:e02_f106_c06_T1_T2_Discharge}.

\begin{longtable}{|p{28mm}|p{22mm}|p{22mm}|p{22mm}|p{22mm}|p{22mm}|}
\caption{Modelled and calculated discharge through the weir for sub-model T1 and T2} 
\label{tab:e02_f106_c06_universal_weir_T1_T2} \\ 
\hline
\STRUT
      {Time [yyyy/MM/dd HH:mm]} 
	 &  {Q-\dflowfm [m\textsuperscript{3}/s]} 
	 &  {Q-\sobekthree [m\textsuperscript{3}/s]}
	 &  {Q-Analytical [m\textsuperscript{3}/s]}
	 &  {Q-\dflowfm \textbf{-} Q-\sobekthree [m\textsuperscript{3}/s]}
	 &  {Q-\dflowfm \textbf{-} Q-Analytical [m\textsuperscript{3}/s]}\\ 
[1ex]\hline \endhead
\hline
\endfoot
%
\STRUT						   
01/01/1996 00:00    & 0.000      & 0.000     & 0.000     & 0.000   & 0.000      \\ 
01/01/1996 06:00    & 5.959      & 5.959     & 5.959     & 0.000   & 0.000      \\
01/01/1996 12:00    & 7.846      & 7.847     & 7.846     & 0.000   & 0.000      \\
01/01/1996 18:00    & 8.936      & 8.936     & 8.936     & 0.000   & 0.000      \\
02/01/1996 00:00    & 9.575      & 9.577     & 9.575     & 0.000   & 0.000      \\
02/01/1996 06:00    & 10.177     & 10.177    & 10.177    & 0.000   & 0.000      \\
02/01/1996 12:00    & 10.177     & 10.177    & 10.177    & 0.000   & 0.000      \\
02/01/1996 18:00    & 10.177     & 10.177    & 10.177    & 0.000   & 0.000      \\
03/01/1996 00:00    & 0.000      & 0.000     & 0.000     & 0.000   & 0.000      \\
03/01/1996 06:00    & -10.177    & -10.177   & -10.177   & 0.000   & 0.000      \\
03/01/1996 12:00    & -10.177    & -10.177   & -10.177   & 0.000   & 0.000      \\
03/01/1996 18:00    & -10.177    & -10.177   & -10.177   & 0.000   & 0.000      \\
04/01/1996 00:00    & -9.575     & -9.577    & -9.575    & 0.000   & 0.000      \\
04/01/1996 06:00    & -8.936     & -8.936    & -8.936    & 0.000   & 0.000      \\
04/01/1996 12:00    & -7.846     & -7.847    & -7.846    & 0.000   & 0.000      \\
04/01/1996 18:00    & -5.959     & -5.959    & -5.959    & 0.000   & 0.000      \\ [1ex]                                                                                                                       
\end{longtable}	                                        

\begin{figure}[H]
    \centering
    \includegraphics*[width=0.95\textwidth]{figures/e02_f106_c06_T1_T2_Discharge.JPG}
    \caption{Timeseries of modelled and calculated discharge through the weir of sub-model T1 and T2}
	\label{fig:e02_f106_c06_T1_T2_Discharge}
\end{figure}

%---------------------------------------------------------------------------------
\section*{Analysis of results}
The \dflowfm modelled discharges through the universal weir are analysed against the results of the analytical solution and \sobekthree model. 
The \dflowfm  and \sobekthree results are very similar as the analytical solution that was computed using the structure equation for the universal weir.


%---------------------------------------------------------------------------------

\printrefsegment

